\documentclass[letterpaper,11pt]{moderncv}
\moderncvstyle{banking}
\usepackage{xpatch}
\moderncvcolor{black}
\AtBeginDocument{\definecolor{color2}{rgb}{0,0,0}}
\usepackage[letterpaper, margin=0.6in]{geometry}
\usepackage{xcolor}
\usepackage{graphicx}
\usepackage{fontspec}
\setmainfont{Times New Roman}
\pagenumbering{arabic}
\usepackage[unicode]{hyperref}
\sloppy
\usepackage{enumitem}

\usepackage[english]{babel}
\usepackage[autostyle, english = american]{csquotes}
\MakeOuterQuote{"}

%==================== PERSONAL INFO ====================%
\name{\textcolor{black}{Saeyoung}}{Kim}
\address{Cambridge, MA, USA}
\phone[mobile]{+1 (351) 322 5510}
\email{saeyoung\_kim@uml.edu}
\social[linkedin]{saeyoung-macx-kim}
\social[github]{PrimaryAnomaly}

\begin{document}

\makecvtitle
\vspace{-3em}

%==================== SUMMARY ====================%
\section{Professional Summary}
Systems test and integration engineer with hands-on experience developing test plans, building test infrastructure, and validating complex electromechanical systems through structured qualification campaigns. Led cross-functional test efforts for military satellite programs, including requirements development, system-level test execution, root cause analysis, and compilation of acceptance evidence across engineering, qualification, and flight model builds. Developed Python-based data analysis tools for automated anomaly detection and statistical process monitoring in manufacturing. Background in sensor systems, embedded hardware/software integration, and translating safety-critical test requirements into validated test procedures.

\section{Technical Skills}
\begin{minipage}[t]{0.48\textwidth}
\cvitem{Test Engineering}{Systems-Level V\&V, Requirements Development, Test Plan and Procedure Authoring, Acceptance Testing, Environmental Qualification}
\cvitem{Electronics}{Circuit Design, PCB Layout (KiCad), Soldering, PCBA Rework, Schematics and Datasheet Analysis}
\cvitem{Programming}{Python, C/C++, MATLAB, Bash, Git, Linux, Windows}
\end{minipage}%
\hfill
\begin{minipage}[t]{0.48\textwidth}
\cvitem{Data and Analysis}{Statistical Process Control, Anomaly Detection, Root Cause Analysis, Data-Driven Safety Evidence, Process Monitoring}
\cvitem{Instruments}{Oscilloscopes, DMMs, Function Generators, Power Supplies, Custom Test Fixtures}
\cvitem{Integration}{Sensor Fusion, Embedded Firmware, UART, SPI, I2C, Data Acquisition Systems}
\cvitem{CAD}{SolidWorks, Fusion360, AutoCAD, 3D Printing}
\end{minipage}

%==================== EXPERIENCE ====================%
\section{Professional Experience}

\cventry{2025.07--Present}{\textbf{Postdoctoral Associate} | Statistical Process Monitoring and Data-Driven Diagnostics}{University of Massachusetts Lowell}{Lowell, MA, USA}{}{
\begin{itemize}[leftmargin=1cm, itemsep=0pt, parsep=0pt]
    \item Built Python-based diagnostic software suite for real-time manufacturing process monitoring, replacing commercial platforms
    \item Developed automated anomaly detection and statistical quality prediction pipelines across 100+ production batches
    \item Compiled and analyzed process data to identify systematic deviations and build evidence-based quality assessments
    \item Designed data-driven test strategies with automated error flagging for multi-day batch manufacturing cycles
    \item Translated cross-functional stakeholder requirements into standardized test procedures and technical documentation
\end{itemize}}

\cventry{2024.02--2025.05}{\textbf{Systems Integration Engineer} | Satellite Test and Validation}{Hanwha Systems}{Seoul, South Korea}{}{
\begin{itemize}[leftmargin=1cm, itemsep=0pt, parsep=0pt]
    \item Led system-level test campaigns for military SAR satellite qualification across EM, QM, and FM builds, developing test plans, validation strategies, and acceptance criteria
    \item Designed Electrical Ground Support Equipment (EGSE) architecture and built test stands for satellite subsystem validation
    \item Compiled requirements, analysis, and test evidence to support formal subsystem and system-level acceptance reviews
    \item Performed structured root cause analysis of electromechanical failures during integration, driving corrective actions to closure
    \item Provided design-for-test and design-for-safety feedback to hardware teams based on failure trends across build campaigns
    \item Coordinated cross-functional technical efforts across electrical, mechanical, software, and program teams throughout qualification
    \item Authored and maintained controlled test procedures, system problem reports, and technical documentation for engineering and customer review
    \item Trained engineers and technicians on test equipment operation, debug processes, and structured troubleshooting methods
    \item Supported environmental testing campaigns including thermal vacuum and vibration qualification
\end{itemize}}

%==================== EDUCATION ====================%
\section{Education}
\cventry{2017.03--2024.02}{\textbf{PhD in Electrical Engineering} | Embedded Sensor Systems and Hardware/Software Integration}{Korea University}{Seoul, South Korea}{}{
\begin{itemize}[leftmargin=1cm, itemsep=0pt, parsep=0pt]
    \item Secured and managed \$220K research grant for CubeSat biological payload hardware development
    \item Designed multi-modal sensor suite integrating electrochemical and optical sensing through full development cycle: requirements, circuit design, PCB layout, fabrication, firmware, and system-level qualification
    \item Developed MCU firmware in C/C++ for real-time multi-channel data acquisition over UART, SPI, and I2C interfaces
    \item Built custom test stands and data acquisition setups using oscilloscopes, DMMs, function generators, and power supplies for sensor characterization and validation
    \item Performed iterative environmental and reliability testing to identify failure modes and validate design changes across prototype revisions
    \item Diagnosed hardware/firmware interaction failures and implemented corrective actions through multiple revision cycles
    \item Published 3 first-author and 10+ co-authored journal papers on sensors, test systems, and space payload hardware
\end{itemize}}

\cventry{2013.03--2017.02}{\textbf{BSc in Mechanical Engineering}}{Konkuk University}{Seoul, South Korea}{}{
\begin{itemize}[leftmargin=1cm, itemsep=0pt, parsep=0pt]
    \item Conducted research in electrochemical sensor microfabrication, cleanroom processing, and hardware prototyping
    \item Gained hands-on experience in CAD modeling, mechanical assembly, PCB prototyping, and precision measurement
\end{itemize}}

%==================== PUBLICATIONS ====================%
\section{Selected Publications}
\begin{sloppypar}
\begin{itemize}[leftmargin=1cm, itemsep=0pt, parsep=0pt]
    \item Jinhwee Kil, \textbf{\underline{Saeyoung Kim}}, James Jungho Pak, "Investigation of thermal management strategies for biological CubeSat payloads," \textit{Acta Astronautica}, 2024, Art. no. 228.
    \item \textbf{\underline{Kim, S.}}; Park, S.; Pak, J.J., "Multi-Modal Multi-Array Electrochemical and Optical Sensor Suite for a Biological CubeSat Payload," \textit{MDPI Sensors}, 2024, 24, Art. no. 265.
    \item \textbf{\underline{Saeyoung Kim}}, Ji-Hoon Han, Beelee Chua, James Jungho Pak, "A pH Sensing Pipette for Cross-Contamination Prevention in Industrial Fermentation," \textit{IEEE Trans. Industrial Electronics}, 2022, 69(7), 7461-7469.
\end{itemize}
\end{sloppypar}

\vfill
\centering{References available upon request.}

\end{document}
